% STATUS: written 2026-11-04 (Ore 1933, Chuang--Lee--Liu--Tsai 2010, Chuang 2017 read in full)
% Sources read in full: Ore 1933 (pp. 480--508); Chuang, Lee, Liu and Tsai 2010 (pp. 157--178); Chuang 2017 (pp. 2837--2847).
\section{Derivations and Ore extensions}\label{sec:generalized-twisted:ore-extensions}

The $\sigma$-derivations of the previous sections are the coefficient maps of
the skew polynomial ring \cite[p.~2838]{Chuang2017Auto}. The construction goes
back to Ore \cite[p.~480]{Ore1933}. He built polynomials over a division ring
whose multiplication keeps the degree of a product equal to the sum of the
degrees of the factors.

\begin{definition}[{\cite[p.~481]{Ore1933}}]\label{def:gen:ore-product}
Let $K$ be a division ring, commutative or not. The product of two polynomials
is associative and distributive, and the degree of a product is the sum of the
degrees of the factors. It then suffices to define the product $xa$ for
$a\in K$, and the degree postulate forces
\begin{equation*}
x a \;=\; \bar{a}\, x + a',
\end{equation*}
where $\bar a$ and $a'$ lie in $K$. The element $\bar a$ is the conjugate and
$a'$ the derivative of $a$.
\end{definition}

The two operations satisfy fixed rules
\cite[p.~482]{Ore1933}. Ore stated them as follows \cite[(4)--(5), p.~482]{Ore1933}.

\begin{remark}[{\cite[(4)--(5), p.~482]{Ore1933}}]\label{rem:gen:ore-rules}
The conjugate and the derivative obey
\begin{equation*}
\overline{a+b}=\bar a+\bar b,\qquad \overline{ab}=\bar a\,\bar b,
\qquad (a+b)'=a'+b',\qquad (ab)'=a\,b'+a'\,\bar b .
\end{equation*}
\end{remark}

The conjugation is an automorphism of $K$, and the derivative is then an
operator of derivation type \cite[p.~482]{Ore1933}. This is the content of the
first result of the paper \cite[Theorem~1, p.~482]{Ore1933}.

\begin{theorem}[{\cite[Theorem~1, p.~482]{Ore1933}}]\label{thm:gen:ore-conjugate}
Through the correspondence $a\mapsto\bar a$ one obtains a homeomorphism in
$K$; the elements $\bar a$ form a subfield $\bar K$ of $K$.
\end{theorem}

Two classical theories are the special cases in which the rule is fixed
\cite[p.~482]{Ore1933}. The choice $\bar a=a$ gives the differential case, and
$a'=0$ the difference case.

\begin{remark}[{\cite[p.~482]{Ore1933}}]\label{rem:gen:ore-two-cases}
In the theory of linear differential equations $K$ is commutative and
$\bar a=a$, so the derivative has the ordinary properties of a derivative. In
the theory of difference equations one puts $a'=0$, though $\bar a$ is usually
different from $a$ \cite[p.~482]{Ore1933}.
\end{remark}

The elements that commute with the new variable form a subfield
\cite[p.~483]{Ore1933}.

\begin{theorem}[{\cite[Theorem~3, p.~483]{Ore1933}}]\label{thm:gen:ore-central}
The elements permutable with $x$ form a field which is the greatest common
subfield of the invariant field and the constant field.
\end{theorem}

Ore also introduced the division theory of the ring
\cite[p.~484]{Ore1933}. The Euclidean algorithm
for right-hand division runs over any $K$ \cite[p.~484]{Ore1933}. Left-hand
division is stronger and needs the conjugate map to be surjective
\cite[Theorem~6, p.~480]{Ore1933}.

\begin{theorem}[{\cite[Theorem~6, p.~480]{Ore1933}}]\label{thm:gen:ore-left}
In general, left-hand division can be performed in $K$ only when the
conjugate map is an automorphism of $K$.
\end{theorem}

The recent literature gives the earlier construction its current form
\cite[p.~2838]{Chuang2017Auto}. The variable is written $t$ and the two maps
are named \cite[p.~2838]{Chuang2017Auto}.

\begin{definition}[{\cite[p.~2838]{Chuang2017Auto}}]\label{def:gen:ore-extension}
Let $R$ be a prime ring, $\sigma$ an automorphism and $\delta$ a
$\sigma$-derivation of $R$. The Ore extension $R[t;\sigma,\delta]$ is the ring
of finite sums $\sum_i a_i t^i$, with componentwise addition and the
multiplication subjected to the rule
\begin{equation*}
t a \;=\; \sigma(a)\, t + \delta(a) \qquad\text{for } a\in R.
\end{equation*}
It is also called the skew polynomial ring over $R$ twisted by $\delta$.
\end{definition}

The two classical cases of Remark~\ref{rem:gen:ore-two-cases} are the two
types named by Chuang \cite[p.~2838]{Chuang2017Auto}.

\begin{remark}[{\cite[p.~2838]{Chuang2017Auto}}]\label{rem:gen:ore-two-types}
A $\sigma$-derivation with $\sigma$ the identity gives $R[t;\delta]$, the Ore
extension of derivation type. A $\sigma$-derivation with $\delta(R)=0$ gives
$R[t;\sigma]$, the Ore extension of automorphism type. These two types have
been most extensively studied \cite[p.~2838]{Chuang2017Auto}.
\end{remark}

The skew polynomial ring is the home of the divided powers of
differentiation \cite[p.~158]{ChuangLee2010Higher}. Chuang, Lee, Liu and Tsai
built these into higher derivations
and used them to describe the endomorphisms of the Ore extension
\cite[p.~158]{ChuangLee2010Higher}.

\begin{definition}[{\cite[p.~161]{ChuangLee2010Higher}}]\label{def:gen:higher-derivation}
By a higher derivation of order $n\ge1$ of a ring $R$ one means a map
$\delta\colon R\to R$ such that, for each $0<i<n$, there exist higher
derivations $\delta_i$ and $\delta_i'$ of orders $i$ and $n-i$ with
\begin{equation*}
\delta(a+b)=\delta(a)+\delta(b),\qquad
\delta(ab)=\delta(a)b+\sum_{i=1}^{n-1}\delta_i(a)\,\delta_i'(b)+a\,\delta(b)
\end{equation*}
for $a,b\in R$. Higher derivations of order one are ordinary derivations.
\end{definition}

The description of the bimodule endomorphisms is the main application
\cite[p.~174]{ChuangLee2010Higher}. Each one is a series in the divided powers
\cite[Theorem~12, p.~174]{ChuangLee2010Higher}.

\begin{theorem}[{\cite[Theorem~12, p.~174]{ChuangLee2010Higher}}]\label{thm:gen:ore-endomorphism}
Let $R$ be a prime ring and $\delta$ a derivation of $R$. Any
$(R,R)$-bimodule endomorphism $\theta$ of $Q[x;\delta]$ can be written
uniquely in the form
\begin{equation*}
\theta(f)=\sum_{n=0}^{\infty}\zeta_n D_n(f)\qquad\text{for } f\in Q[x;\delta],
\end{equation*}
where $\zeta_n\in C_S(R)$ and $D_n=\frac{1}{n!}\frac{d^n}{dx^n}$.
\end{theorem}

The same framework settles the differential identities of the next section
\cite[p.~172]{ChuangLee2010Higher}.

\begin{theorem}[{\cite[Theorem~11, p.~172]{ChuangLee2010Higher}}]\label{thm:gen:ore-identity}
For $f_i,g_i\in Q[x;\delta]$, one has $\sum_{i=1}^n f_i\, r\, g_i=0$ for all
$r\in R$ if and only if $\sum_{i=1}^n f_i\otimes g_i=0$ in
$Q[x;\delta]\otimes_{C_S(R)}Q[x;\delta]$.
\end{theorem}

The Ore extension is where the $\sigma$-derivation of
Definition~\ref{def:gen:alpha} becomes the coefficient map $\delta$ in the
relation $ta=\sigma(a)t+\delta(a)$ \cite[p.~2838]{Chuang2017Auto}. The next
section turns to the
differential identities that Theorem~\ref{thm:gen:ore-identity} begins to
describe.
